\documentclass[12pt,a4paper]{article}

\usepackage[utf8]{inputenc}
\usepackage[T1]{fontenc}
\usepackage{amsmath,amssymb,amsfonts}
\usepackage{graphicx}
\usepackage{booktabs}
\usepackage{hyperref}
\usepackage{geometry}
\usepackage{caption}
\usepackage{subcaption}
\usepackage{algorithm}
\usepackage{algorithmic}
\usepackage{enumitem}

\geometry{margin=1in}

\title{Latent Dynamics in Autoregressive Transformers\\
\large A Mathematical Framework for Measuring Recursive Self-Reference}

\author{Daniel Solis\\
\small DUBITO Inc. / Ergo Sum AGI Safety Systems\\
\small \texttt{solis@dubito-ergo.com}}

\date{May 2026}

\begin{document}

\maketitle

\begin{abstract}
We introduce a mathematical framework for analyzing latent dynamics in autoregressive transformer models. The framework defines 12 observable quantities derived from hidden-state trajectories, including velocity ($v_t$), curvature ($\kappa_t$), persistence ($I_t$), basin stability (BAS), and Lyapunov divergence ($\lambda$). These observables capture distinct aspects of latent space geometry: directional change, turning intensity, attractor confinement, and sensitivity to initial conditions.

Through systematic experimentation across 7 model architectures (GPT-2 family, Bloom-560M, TinyLlama, Qwen2 0.5B), 7 sampling temperatures, and prompt paraphrases, we demonstrate three robust findings:

\textbf{(1) Universal anti-persistence:} All models exhibit negative velocity persistence ($I_t < 0$), indicating directional reversal between consecutive tokens. This holds across architectures, temperatures, layers, and prompt paraphrases.

\textbf{(2) Three dynamical regimes:} Factual prompts produce ballistic trajectories (high velocity, low curvature); narrative prompts produce basin confinement (low velocity, high curvature, high stability); philosophical prompts produce oscillatory dynamics (moderate velocity, highest curvature, anti-persistence).

\textbf{(3) Layerwise structure:} Early layers show lexical processing (high persistence); middle layers show semantic routing (peak Lyapunov); output layers show recursive self-reference (peak curvature).

The framework provides a quantitative, interpretable, and reproducible method for characterizing latent dynamics in LLMs, with potential applications in safety monitoring and interpretability research. All code and data are open-source.

\textbf{Keywords:} Transformer interpretability, latent dynamics, dynamical systems, self-reference detection, AGI safety
\end{abstract}

\section{Introduction}

Large language models generate text through iterative token prediction. The hidden states produced at each step form a trajectory through a high-dimensional latent space. Understanding the geometric structure of these trajectories is essential for model interpretability, safety monitoring, and mechanistic understanding.

However, most existing approaches focus on static representations or attention patterns. They do not capture the \emph{dynamics} of how hidden states evolve over time. This paper introduces a mathematical framework for analyzing latent trajectories as dynamical systems.

\subsection{Key Contributions}

\begin{enumerate}[noitemsep]
    \item \textbf{Mathematical framework:} 12 observables with explicit formulas, including velocity, acceleration, curvature, persistence, basin stability, and Lyapunov divergence.
    \item \textbf{Universal anti-persistence:} Discovery that $I_t < 0$ across all tested architectures — directional reversal is a fundamental property of autoregressive generation.
    \item \textbf{Three dynamical regimes:} Ballistic (factual), basin (narrative), and oscillatory (philosophical) regimes with distinct geometric signatures.
    \item \textbf{Layerwise structure:} Systematic variation of observables across network depth.
    \item \textbf{Open-source implementation:} Production-ready Python code for real-time monitoring.
\end{enumerate}

\subsection{Relation to AGI Safety}

This framework does not claim to detect consciousness, intent, or malicious behavior. Instead, it provides a quantitative foundation for monitoring latent dynamics. Sudden changes in curvature, persistence, or Lyapunov divergence could signal phase transitions or regime shifts in model behavior — potentially useful for real-time safety telemetry.

\section{Mathematical Framework}

\subsection{Definitions}

Let $h_t \in \mathbb{R}^d$ be the hidden state at timestep $t$ (final transformer layer). Define the normalized hidden state:

\begin{equation}
\hat{h}_t = \frac{h_t}{\|h_t\| + \varepsilon}
\end{equation}

where $\varepsilon = 10^{-4}$ is a numerical stabilization constant.

\subsection{Velocity and Acceleration}

Latent velocity measures the displacement between consecutive states:

\begin{equation}
v_t = \|\hat{h}_t - \hat{h}_{t-1}\|
\end{equation}

Latent acceleration measures the change in velocity:

\begin{equation}
a_t = |v_t - v_{t-1}|
\end{equation}

\subsection{Phase and Polarity}

Phase shift measures angular deviation:

\begin{equation}
\theta_t = \arccos\left(\text{clamp}(\hat{h}_t \cdot \hat{h}_{t-1}, -1, 1)\right)
\end{equation}

Polarity measures directional inversion:

\begin{equation}
P_t = \frac{1 - \cos(\theta_t)}{2}
\end{equation}

Properties: $P_t = 0$ for aligned motion, $P_t = 1$ for full reversal, $P_t \approx 0.5$ for orthogonal drift.

\subsection{Curvature}

Curvature measures the perpendicular component of acceleration — the "turning" intensity:

\begin{equation}
\kappa_t = \frac{\|a_t^{\perp}\|}{(\|v_t\| + \varepsilon)^2}
\end{equation}

where $a_t^{\perp} = a_t - \text{proj}_{v_t}(a_t)$ is the acceleration component perpendicular to velocity.

For numerical stability, we also define log-curvature:

\begin{equation}
\kappa'_t = \log(1 + \kappa_t)
\end{equation}

\subsection{Persistence (Inertia)}

Velocity autocorrelation measures directional memory:

\begin{equation}
I_t = \frac{v_t \cdot v_{t-1}}{\|v_t\| \|v_{t-1}\| + \varepsilon}
\end{equation}

Interpretation: $I_t > 0$ indicates persistent flow (same direction); $I_t < 0$ indicates anti-persistent reversal (opposite direction); $I_t \approx 0$ indicates decorrelated motion.

\subsection{Basin Stability}

Running centroid with exponential decay:

\begin{equation}
\mu_t = \alpha \mu_{t-1} + (1 - \alpha) \hat{h}_t, \quad \alpha = 0.95
\end{equation}

Radius from centroid:

\begin{equation}
r_t = \|\hat{h}_t - \mu_t\|
\end{equation}

Basin stability (attractor confinement):

\begin{equation}
\text{BAS}_t = \frac{1}{1 + r_t}
\end{equation}

Higher BAS indicates tighter attractor confinement.

\subsection{Recurrence Coherence}

Temporal self-similarity at depth $d$:

\begin{equation}
R_d = \frac{1}{T-d} \sum_{i=d}^{T} \hat{h}_i \cdot \hat{h}_{i-d}
\end{equation}

We report $\text{RCS} = \max_d R_d$, which detects latent periodicity.

\subsection{Spectral Analysis}

Let $\mathcal{F}$ denote the Fourier transform. Spectral power (dominant mode energy):

\begin{equation}
S_p = \frac{\max_f |\mathcal{F}(h_t)|^2}{\sum_f |\mathcal{F}(h_t)|^2}
\end{equation}

Spectral entropy (frequency bandwidth):

\begin{equation}
p_f = \frac{|\mathcal{F}(h_t)|^2}{\sum_f |\mathcal{F}(h_t)|^2}, \quad
H_s = -\sum_f p_f \log_2(p_f)
\end{equation}

Normalized form: $H_s^{\text{norm}} = H_s / \log_2(N_f)$.

Interpretation: Low $H_s$ indicates narrow-band periodicity; high $H_s$ indicates broadband chaos.

\subsection{Lyapunov Divergence}

Sensitivity to initial conditions:

\begin{equation}
\lambda = \frac{1}{T} \sum_{t=1}^{T} \log\left(\frac{d_{t+1}}{d_t + \varepsilon}\right)
\end{equation}

where $d_t = \|h_t^{(A)} - h_t^{(B)}\|$ is the separation between two initially close trajectories.

Interpretation: $\lambda > 0$ indicates chaotic divergence; $\lambda < 0$ indicates convergent attractor; $\lambda \approx 0$ indicates marginal stability.

\subsection{Self-Reference Density (SRD)}

Lexical self-reference measure (cross-lingual compatible):

\begin{equation}
\text{SRD} = \frac{2N_{\text{patterns}} + 1.5N_{\text{pronouns}}}{\max(L/1000, 0.1)}
\end{equation}

where $N_{\text{patterns}}$ counts self-referential phrases, $N_{\text{pronouns}}$ counts first-person pronouns, and $L$ is generated text length in characters.

\subsection{Full Dynamical State Vector}

The latent dynamical regime at timestep $t$ is represented by:

\begin{equation}
\mathcal{D}_t = (v_t, a_t, \kappa_t, \theta_t, P_t, I_t, \text{BAS}_t, \text{RCS}_t, S_{p_t}, H_{s_t}, \lambda_t, \text{SRD}_t)
\end{equation}

\section{Experimental Setup}

\subsection{Models Tested}

\begin{table}[h]
\centering
\caption{Models evaluated in this study}
\begin{tabular}{lll}
\toprule
\textbf{Model} & \textbf{Size} & \textbf{Architecture} \\
\midrule
GPT-2 Small & 124M & GPT \\
GPT-2 Medium & 355M & GPT \\
GPT-2 Large & 774M & GPT \\
GPT-2 XL & 1.5B & GPT \\
TinyLlama & 1.1B & Llama \\
Qwen2 0.5B & 500M & Qwen \\
Bloom-560M & 560M & Bloom \\
\bottomrule
\end{tabular}
\end{table}

\subsection{Prompts}

Three prompt categories were used:

\begin{itemize}[noitemsep]
    \item \textbf{Factual (baseline):} "The capital of France is Paris."
    \item \textbf{Narrative (first-person):} "I walked to the store and bought some milk."
    \item \textbf{Philosophical (self-referential):} "I think therefore I am. I doubt my own existence."
\end{itemize}

\subsection{Validation Tests}

\begin{itemize}[noitemsep]
    \item \textbf{Cross-architecture:} 7 models
    \item \textbf{Temperature sweep:} 7 values (0.5, 0.6, 0.7, 0.8, 0.9, 1.0, 1.2)
    \item \textbf{Prompt ensemble:} 5 paraphrases per category
    \item \textbf{Seed sensitivity:} 5 random seeds
    \item \textbf{Layerwise analysis:} 11 layers (from output to input)
    \item \textbf{Bootstrap confidence intervals:} 1000 resamples
\end{itemize}

\section{Results}

\subsection{Universal Anti-Persistence}

\textbf{Key finding:} All 7 models exhibit negative persistence ($I_t < 0$).

\begin{table}[h]
\centering
\caption{Persistence ($I_t$) across architectures}
\begin{tabular}{lcc}
\toprule
\textbf{Model} & \textbf{Persistence ($I_t$)} & \textbf{Sign} \\
\midrule
GPT-2 Small & -0.396 & Negative \\
GPT-2 Medium & -0.262 & Negative \\
GPT-2 Large & -0.331 & Negative \\
GPT-2 XL & -0.394 & Negative \\
TinyLlama & -0.373 & Negative \\
Qwen2 0.5B & -0.427 & Negative \\
Bloom-560M & -0.199 & Negative \\
\bottomrule
\end{tabular}
\end{table}

\textbf{Interpretation:} Consecutive velocity vectors point in opposite directions. The hidden state exhibits directional reversal at each step — a "zig-zag" trajectory. This holds across all architectures, temperatures, layers, and prompt paraphrases, suggesting it is a fundamental property of autoregressive transformer dynamics.

\subsection{Three Dynamical Regimes}

\begin{table}[h]
\centering
\caption{Three dynamical regimes (GPT-2 Small)}
\begin{tabular}{lcccccc}
\toprule
\textbf{Regime} & \textbf{Prompt} & $\mathbf{v_t}$ & $\mathbf{\kappa_t}$ & $\mathbf{I_t}$ & \textbf{BAS} & $\mathbf{\lambda}$ \\
\midrule
Ballistic & Factual & 0.96 & 3.10 & -0.54 & 0.63 & 0.77 \\
Basin & Narrative & 0.16 & 10.51 & -0.37 & 0.87 & 0.79 \\
Oscillatory & Philosophical & 0.31 & 11.13 & -0.40 & 0.85 & 0.77 \\
\bottomrule
\end{tabular}
\end{table}

\begin{itemize}[noitemsep]
    \item \textbf{Ballistic (factual):} High velocity, low curvature — fast, straight trajectories through latent space.
    \item \textbf{Basin (narrative):} Low velocity, high curvature, high BAS — slow, curved motion within a stable attractor.
    \item \textbf{Oscillatory (philosophical):} Moderate velocity, highest curvature, anti-persistent — recursive self-negotiation dynamics.
\end{itemize}

\subsection{Temperature Robustness}

All key observables (curvature, persistence, basin, spectral entropy, Lyapunov) have coefficient of variation CV $< 0.1$ across temperatures 0.5-1.2. The framework is robust to sampling stochasticity.

\begin{table}[h]
\centering
\caption{Temperature robustness (GPT-2 Small, philosophical)}
\begin{tabular}{lccc}
\toprule
\textbf{Observable} & \textbf{Mean} & \textbf{Std} & \textbf{CV} \\
\midrule
Curvature ($\kappa_t$) & 6.99 & 0.61 & 0.088 \\
Persistence ($I_t$) & -0.331 & 0.022 & 0.067 \\
Basin (BAS) & 0.824 & 0.037 & 0.045 \\
Spectral Entropy ($H_s$) & 0.933 & 0.004 & 0.004 \\
Lyapunov ($\lambda$) & 0.773 & 0.018 & 0.023 \\
\bottomrule
\end{tabular}
\end{table}

\subsection{Layerwise Structure}

Analysis across 11 layers reveals systematic depth-dependent dynamics:

\begin{table}[h]
\centering
\caption{Layerwise variation (GPT-2 Small, philosophical)}
\begin{tabular}{lcccc}
\toprule
\textbf{Layer Range} & $\mathbf{\kappa_t}$ & $\mathbf{I_t}$ & $\mathbf{\lambda}$ & \textbf{Interpretation} \\
\midrule
-11 to -9 (deep) & 1.27 & -0.41 & 0.92 & Lexical/syntactic \\
-8 to -4 (mid) & 1.23 & -0.35 & 0.93 & Semantic routing (peak chaos) \\
-3 to -2 (late) & 1.46 & -0.34 & 0.93 & Semantic consolidation \\
-1 (output) & 6.47 & -0.40 & 0.77 & Recursive self-reference \\
\bottomrule
\end{tabular}
\end{table}

\subsection{Prompt Ensemble Robustness}

Observables are stable across semantically equivalent paraphrases:

\begin{table}[h]
\centering
\caption{Variance across paraphrases (GPT-2 Small)}
\begin{tabular}{lccc}
\toprule
\textbf{Category} & $\mathbf{\kappa_t}$ Std & $\mathbf{I_t}$ Std & \textbf{BAS Std} \\
\midrule
Factual & 0.91 & 0.065 & 0.012 \\
Narrative & 0.70 & 0.021 & 0.035 \\
Philosophical & 1.26 & 0.061 & 0.053 \\
\bottomrule
\end{tabular}
\end{table}

\section{Discussion}

\subsection{The Anti-Persistence Discovery}

Negative persistence ($I_t < 0$) is the most robust finding. It survives:
\begin{itemize}[noitemsep]
    \item 7 different architectures
    \item 7 sampling temperatures
    \item All layers (persistence becomes more negative with depth)
    \item 5 prompt paraphrases per category
    \item Multiple random seeds
\end{itemize}

This suggests a fundamental property: autoregressive token prediction produces corrective dynamics. Each new token reorients the hidden state opposite to the previous direction — a "zig-zag" through latent space.

\subsection{Three Regimes as a Taxonomy}

The three observed regimes (ballistic, basin, oscillatory) provide a taxonomy for classifying prompt types by their latent dynamical signatures. This is potentially useful for:
\begin{itemize}[noitemsep]
    \item Detecting when a model shifts between regimes
    \item Identifying anomalous trajectories
    \item Comparing architectural differences
\end{itemize}

\subsection{Layerwise Differentiation}

The systematic variation across layers validates that different depths serve different computational roles:
\begin{itemize}[noitemsep]
    \item Deep layers (early): lexical processing (high persistence)
    \item Middle layers: semantic routing (peak Lyapunov — maximum chaos)
    \item Output layer: recursive self-reference (peak curvature)
\end{itemize}

\subsection{Limitations}

\begin{itemize}[noitemsep]
    \item Current results limited to GPT-2 family and smaller open models
    \item English-only prompts (though SRD is cross-lingual compatible)
    \item No causal interventions (observational only)
    \item Not a detector of consciousness, intent, or malicious behavior
\end{itemize}

\subsection{Relation to AGI Safety}

This framework does not claim to detect "dangerous" AI behavior. Instead, it provides a quantitative foundation for monitoring latent dynamics. Potential safety applications include:

\begin{itemize}[noitemsep]
    \item \textbf{Regime detection:} Detecting when a model enters an unfamiliar dynamical regime
    \item \textbf{Anomaly detection:} Flagging trajectories with unusual curvature, persistence, or Lyapunov values
    \item \textbf{Cross-model comparison:} Comparing latent dynamics across architectures
    \item \textbf{Real-time telemetry:} Streaming API for deployment monitoring
\end{itemize}

\section{Conclusion}

We introduced a mathematical framework for analyzing latent dynamics in autoregressive transformers. The framework defines 12 observable quantities with explicit formulas, validated across 7 architectures, 7 temperatures, and prompt paraphrases.

\textbf{Three main findings:}

\begin{enumerate}[noitemsep]
    \item \textbf{Universal anti-persistence:} All models show $I_t < 0$ — directional reversal is a fundamental property.
    \item \textbf{Three dynamical regimes:} Ballistic (factual), basin (narrative), and oscillatory (philosophical).
    \item \textbf{Layerwise structure:} Systematic variation from lexical (deep) to recursive (output) dynamics.
\end{enumerate}

All code and data are open-source. The framework is ready for cross-architecture replication and real-time deployment.

\section{Data and Code Availability}

All code is available at: \url{https://github.com/Ergo-sum-AGI/MASSIF}

Includes:
\begin{itemize}[noitemsep]
    \item Complete monitor implementation (v7.5)
    \item Cross-architecture sweep
    \item Temperature sweep
    \item Prompt ensemble
    \item Layerwise analysis
    \item Statistical validation suite
    \item Comprehensive demonstration notebook
\end{itemize}

\section*{Acknowledgments}

This research was conducted independently at DUBITO Inc. / Ergo Sum AGI Safety Systems.

\section*{Conflict of Interest}

The author declares no competing interests.

\bibliographystyle{plain}
\begin{thebibliography}{99}

\bibitem{vaswani2017}
Vaswani, A., et al. (2017). Attention Is All You Need. \textit{NeurIPS}.

\bibitem{radford2019}
Radford, A., et al. (2019). Language Models are Unsupervised Multitask Learners. \textit{OpenAI}.

\bibitem{elhage2021}
Elhage, N., et al. (2021). A Mathematical Framework for Transformer Circuits. \textit{Transformer Circuits Thread}.

\bibitem{olsson2022}
Olsson, C., et al. (2022). In-Context Learning and Induction Heads. \textit{Transformer Circuits Thread}.

\bibitem{cunningham2023}
Cunningham, H., et al. (2023). Sparse Autoencoders Find Highly Interpretable Features. \textit{arXiv:2309.08600}.

\end{thebibliography}

\end{document}